\documentclass[10pt,a4paper]{amsart}
\usepackage[margin=19mm]{geometry}
\usepackage{amsmath,amssymb,mathtools}
\usepackage[scr=rsfs,cal=euler]{mathalfa}
\usepackage{xfrac}
\usepackage[unicode,hidelinks,bookmarks=false]{hyperref}
\newcommand{\ie}{i.e.\ }
\newcommand{\cf}{cf.\ }
\newcommand{\resp}{respectively\ }
\newcommand{\Subsection}{\subsection}
\providecommand{\nobreakdash}{\nobreak\hbox{-}}
\setlength{\emergencystretch}{2em}
\multlinegap0pt
\usepackage{bm}
\usepackage{bbm}
\usepackage{dsfont}
\usepackage{xspace}
\usepackage{cancel}
\usepackage{mathtools}
\usepackage{amsthm}
\makeatletter
\@ifundefined{theorem}{\newtheorem{theorem}{Theorem}}{}
\makeatother
\title[High-order Lagrange multiplier schemes for general Hamiltoni]{High-order Lagrange multiplier schemes for general Hamiltonian PDEs}
\author{Yonghui Bo, Yushun Wang}
\date{Source: arXiv:2601.12776v1 (CC BY 4.0)}
\begin{document}
\maketitle

\noindent\emph{Source and licence.} High-order Lagrange multiplier schemes for general Hamiltonian PDEs. Version arXiv:2601.12776v1 (CC BY 4.0), distributed under the Creative Commons Attribution 4.0 International licence, as recorded in the arXiv licence metadata for this exact version and shown on the abstract page https://arxiv.org/abs/2601.12776v1. Full rights notice: licenses/arxiv-2601.12776-CC-BY-4.0.txt.\par\smallskip

\noindent\textit{Editorial scope.} Counted unit: the Theorem whose source label is \texttt{Th-2-6} in the article (source0.tex lines 368--373, source label \texttt{Th-2-6}), delivered together with the article's own complete original proof of it (source0.tex lines 374--446, one \texttt{proof} environment of 73 lines). No mathematics is written, repaired or re-derived: statement and proof are the article's own text line for line. Required context retained uncounted: eq-2-10 (lines 185--189); eq-2-11 (lines 191--193); Th-2-4 (lines 235--240); eq-2-17 (lines 255--259); eq-2-18 (lines 262--269); eq-2-28 (lines 350--352). Every definition, display and label of those lines is retained unchanged. Omitted from this package and not redistributed: 1--184, 190--190, 194--234, 241--254, 260--261, 270--349, 353--367, 447--714. Nothing inside a retained excerpt is omitted or altered: every retained body is delivered exactly as the article's lines are written, so no omission receipt of any kind accompanies this package. Citations: the retained excerpts carry no citation command at all, so no bibliography entry of the article is transcribed and no reference is redistributed. Reference adaptation: none was needed and none was made; every reference inside the delivered unit resolves inside it, so no unresolved reference or citation is produced. Printed slips kept literal: no printed word, symbol, hypothesis or number of the counted statement or of its proof was altered. Layout and class: the article's own class is \texttt{siamart} with packages including amsfonts, graphicx, epstopdf, algorithmic; the wrapper is the lane's pinned \texttt{amsart} editorial wrapper at 10pt on A4, which restates the theorem-like environments the retained text needs and adds the article's own macro definitions that the retained text needs (none), with extra wrapper packages (bm, bbm, dsfont, xspace, cancel, mathtools). The delivered numbering is the unit's own. Assets: no retained span contains a figure environment, a graphics-inclusion command, a PDF-page inclusion, a table or a TikZ picture; no image file is copied into this package, the package declares no asset dependency and no image file is redistributed with it. Licence: version arXiv:2601.12776v1 of this article is distributed under the Creative Commons Attribution 4.0 International licence, as recorded in the arXiv licence metadata for this exact version and shown on the abstract page https://arxiv.org/abs/2601.12776v1. A full rights notice is delivered at licenses/arxiv-2601.12776-CC-BY-4.0.txt. Characters: every retained excerpt is pure ASCII, so the delivered engine, XeLaTeX with its default Latin Modern text font, reports no missing character.
	\begin{equation}\label{eq-2-10}
	\begin{aligned}
		z_{i,(m+1)}^n&=z^n+\Delta t\sum_{j=1}^sa_{ij}k_{j,(m+1)}^n,\\ k_{i,(m+1)}^n&=f_1\big(z_{i,(m+1)}^n\big)+f_2\big(z_{i,(m)}^n\big),
		\end{aligned}
	\end{equation}

	\begin{equation}\label{eq-2-11}
		z_{(\Lambda)}^{n+1}=z^n+\Delta t\sum_{i=1}^sb_ik_{i,(\Lambda)}^n.
	\end{equation}

\begin{theorem}\label{Th-2-4}
Given that the underlying RK method is of order $p$, the scheme \eqref{eq-2-10}-\eqref{eq-2-11} possesses an order of $\min\{p,\Lambda\}$, i.e.,
\begin{equation}\label{eq-2-16}
	z_{(\Lambda)}^{n+1}-z(t_{n+1})=\mathcal{O}\big((\Delta t)^{\min\{p,\Lambda\}+1}\big).
\end{equation}
\end{theorem}

	\begin{equation}\label{eq-2-17}
		\begin{aligned}
			z_{i,(m+1)}^n&=z^n+\Delta t\sum_{j=1}^sa_{ij}k_{j,(m+1)}^n,\\ k_{i,(m+1)}^n&=f_1\big(z_{i,(m+1)}^n\big)+f_2\big(z_{i,(m)}^n\big).
		\end{aligned}
	\end{equation}

	\begin{equation}\label{eq-2-18}
		\begin{aligned}
			z_i^n&=z^n+\Delta t\sum_{j=1}^sa_{ij}k_j^n,\\
			k_i^n&=f_1\big(z_i^n\big)+\lambda^nf_2\big(\tilde{z}_i^n\big),\\
			z^{n+1}&=z^n+\Delta t\sum_{i=1}^sb_ik_i^n,\\
			F(z^{n+1})&=F(z^n)+\lambda^n\Delta t\sum_{i=1}^sb_i\big(N'(\tilde{z}_i^n),k_i^n\big).
		\end{aligned}
	\end{equation}

\begin{equation}\label{eq-2-28}
	\big(N(L_2^n+\lambda^nR_2^n)-N(z^n),1\big)-\lambda^n\Delta t\textbf{b}\big(N'(\tilde{\textbf{z}}^n),L_1^n+\lambda^nR_1^n\big)=0.
\end{equation}

\begin{theorem}\label{Th-2-6}
If $\big(N'(z^n),\frac{\delta\mathcal{H}}{\delta z}(z^n)\big)\neq0$, there exists a constant $\Delta t^*$ such that, for all $\Delta t\in(0,\Delta t^*]$, the algebraic equation \eqref{eq-2-28} locally defines a unique solution $\lambda^n=\lambda^n(\Delta t)$ in a neighborhood of $1$. Furthermore, if the underlying RK method is of order $p$, then the scheme \eqref{eq-2-17}-\eqref{eq-2-18} is of order \mbox{min}\{p,$\Lambda$\}, i.e.,
\begin{equation}\label{eq-2-29}
	z^{n+1}-z(t_{n+1})=\mathcal{O}\big((\Delta t)^{\min\{p,\Lambda\}+1}\big).
\end{equation}
\end{theorem}

\begin{proof}
Note that $\tilde{\textbf{z}}^n$, $L_1^n$, $R_1^n$, $L_2^n$ and $R_2^n$ in \eqref{eq-2-28} can be regarded as functions of $\Delta t$. It follows that
\begin{equation}\label{eq-2-30}
	\begin{aligned}
	\lim_{\Delta t\rightarrow 0}\tilde{\textbf{z}}^n&=\textbf{z}^n,\quad\lim_{\Delta t\rightarrow 0}L_1^n=f_1(\textbf{z}^n),\quad \lim_{\Delta t\rightarrow 0}R_1^n=f_2(\textbf{z}^n),\\
    \lim_{\Delta t\rightarrow 0}L_2^n&=z^n,\quad \lim_{\Delta t\rightarrow 0}R_2^n=0.
	\end{aligned}
\end{equation}
From \eqref{eq-2-28}, we introduce
\begin{equation}\label{eq-2-31}
	\lambda^n=\frac{\big(N(L_2^n+\lambda^nR_2^n)-N(z^n),1\big)}{\Delta t\textbf{b}\big(N'(\tilde{\textbf{z}}^n),L_1^n+\lambda^nR_1^n\big)}:=\varphi(\lambda^n,\Delta t).
\end{equation}
Using the L'H\^{o}pital's rule, along with the antisymmetry of $\mathcal{S}$ and \eqref{eq-2-30}, we derive
\begin{equation}\label{eq-2-32}
	\lim_{\Delta t\rightarrow 0}\lambda^n=\lim_{\Delta t\rightarrow 0}\varphi(\lambda^n,\Delta t)=\frac{\big(N'(z^n),f_1(z^n)\big)}{\big(N'(z^n),f_1(z^n)\big)}=1,
\end{equation}	
where $\big(N'(z^n),f_1(z^n)\big)=\big(N'(z^n),\frac{\delta\mathcal{H}}{\delta z}(z^n)\big)\neq0$. Analogous calculations show that
\begin{equation*}
	\lim_{\Delta t\rightarrow 0}\frac{\partial\varphi(\lambda^n,\Delta t)}{\partial\lambda^n}=\frac{0}{\big(N'(z^n),f_1(z^n)\big)}=0.
\end{equation*}	
This equation confirms a contraction property. Therefore, by the fixed-point theorem, for a sufficiently small $\Delta t^*$, \eqref{eq-2-31} possesses a unique solution $\lambda^n=\lambda^n(\Delta t)$ near $1$ for all $\Delta t\in(0,\Delta t^*]$. Subsequently, we can compute
\begin{equation}\label{eq-2-33}
	\lambda^n-1=\frac{\big(N(L_2^n+\lambda^nR_2^n)-N(z^n),1\big)-\Delta t\textbf{b}\big(N'(\tilde{\textbf{z}}^n),L_1^n+\lambda^nR_1^n\big)}{\Delta t\textbf{b}\big(N'(\tilde{\textbf{z}}^n),L_1^n+\lambda^nR_1^n\big)}:=\frac{\mu(\Delta t)}{\nu(\Delta t)}.
\end{equation}
Combining \eqref{eq-2-30}, \eqref{eq-2-32}, and the identity $\sum_{i=1}^s b_i = 1$, we obtain
\begin{equation*}
	\lim_{\Delta t\rightarrow 0}\textbf{b}\big(N'(\tilde{\textbf{z}}^n),L_1^n+\lambda^nR_1^n\big)=\big(N'(z^n),f_1(z^n)\big)\neq0.
\end{equation*}	
This result implies that
\begin{equation*}
	\nu(\Delta t)=\mathcal{O}(\Delta t).
\end{equation*}
Then, we define $z_{(\Lambda)}^{n+1}$ as the solution of
\begin{equation*}
	\begin{aligned}
		\hat{z}_i^n&=z^n+\Delta t\sum_{j=1}^sa_{ij}\hat{k}_j^n,\\
		\hat{k}_i^n&=f_1\big(\hat{z}_i^n\big)+f_2\big(\tilde{z}_i^n\big),\\
		z_{(\Lambda)}^{n+1}&=z^n+\Delta t\sum_{i=1}^sb_i\hat{k}_i^n,
	\end{aligned}
\end{equation*}
where $\tilde{z}_i^n=z_{i,(\Lambda)}^n$ is computed from \eqref{eq-2-17}. According to Theorem \ref{Th-2-4}, if the underlying RK method is of order $p$, it follows that
\begin{equation}\label{eq-2-34}
	\begin{aligned}
		z_{(\Lambda)}^{n+1}-z(t_{n+1})&=\mathcal{O}\big((\Delta t)^{\min\{p,\Lambda\}+1}\big),\\
		F\big(z(t_{n+1})\big)-\big(F(z^n)+\Delta t\textbf{b}\big(N'(\tilde{\textbf{z}}^n),\hat{\textbf{k}}^n\big)\big)&=\mathcal{O}\big((\Delta t)^{\min\{p,\Lambda\}+1}\big),
    \end{aligned}
\end{equation}	
where $\hat{\textbf{k}}^n=\big(\hat{k}_1^n,\cdots,\hat{k}_s^n\big)^T$, and we have $z_{(\Lambda)}^{n+1}=L_2^n+R_2^n$ and $\hat{\textbf{k}}^n=L_1^n+R_1^n$. Using the Taylor expansion and the first equation in \eqref{eq-2-34}, we get
\begin{equation*}
	F\big(z_{(\Lambda)}^{n+1}\big)-F\big(z(t_{n+1})\big)=\mathcal{O}\big((\Delta t)^{\min\{p,\Lambda\}+1}\big).
\end{equation*}
With \eqref{eq-2-32}, the above equations imply that as $\Delta t\rightarrow 0$,
\begin{equation*}
	\begin{aligned}
		\mu(\Delta t)&=\big(N(L_2^n+\lambda^nR_2^n)-N(z^n),1\big)-\Delta t\textbf{b}\big(N'(\tilde{\textbf{z}}^n),L_1^n+\lambda^nR_1^n\big),\\
		&=F\big(z_{(\Lambda)}^{n+1}\big)-F\big(z(t_{n+1})\big)+F\big(z(t_{n+1})\big)-\big(F(z^n)+\Delta t\textbf{b}\big(N'(\tilde{\textbf{z}}^n),\hat{\textbf{k}}^n\big)\big)\\
		&=\mathcal{O}\big((\Delta t)^{\min\{p,\Lambda\}+1}\big).
	\end{aligned}
\end{equation*}	
From \eqref{eq-2-33}, we then have
\begin{equation}\label{eq-2-35}
	\lambda^n-1=\frac{\mu(\Delta t)}{\nu(\Delta t)}=\mathcal{O}\big((\Delta t)^{\min\{p,\Lambda\}}\big).
\end{equation}
Thus, the accuracy of the scheme \eqref{eq-2-17}-\eqref{eq-2-18} can be obtained from
\begin{equation*}
	\begin{aligned}
		z^{n+1}-z(t_{n+1})&=L_2^n+\lambda^nR_2^n-z(t_{n+1})\\
		&=L_2^n+R_2^n-z(t_{n+1})+(\lambda^n-1)R_2^n,\\
		&=z_{(\Lambda)}^{n+1}-z(t_{n+1})+(\lambda^n-1)\Delta t\textbf{b}R_1^n.
	\end{aligned}
\end{equation*}	
In the limit as $\Delta t\rightarrow 0$, the equations \eqref{eq-2-30}, \eqref{eq-2-34} and \eqref{eq-2-35} collectively yield \eqref{eq-2-29}, which completes the proof.
\end{proof}

\end{document}
